%!TEX root=./main.tex

\section{Introduction}

Given a concurrent program, the \emph{estimation problem} asks: how many different behaviors (up to a certain size) does it have?
It is a natural, counting analogue of decision questions such as reachability.
%
In this work, we develop systematic solutions for the estimation problem.
As a first step, we must answer what we mean by a behavior, since 
different model checkers induce different notions of behaviors.
A naive interleaving-based model checker treats every thread schedule as a distinct execution.
In contrast, a checker based on partial-order reduction \cite{GodefroidThesis, spin, dpor, verisoft,pop-dpor,optimal-dpor,nidhugg,trust} 
collapses all schedules within the same equivalence class into a single behavior, 
thereby exploring a significantly smaller space.
Additional heuristics may reduce the space further, e.g., by defining coarser equivalences based on symmetry.

We focus on behaviors induced by Mazurkiewicz's trace equivalence, where two executions are considered equivalent
if one can be obtained from the other by repeated swapping of adjacent commutative operations.
Our choice is guided, on the one hand, by the success of this notion in model checking
\cite{spin, verisoft, dpor,optimal-dpor, trust}
and, on the other, by the recent emergence of many \emph{optimal} exploration algorithms that guarantee that each equivalence class is visited exactly 
once \cite{optimal-dpor, trust, must}.
In this setting, the estimation problem asks:
\begin{quote}
{\bf Exact Estimation problem:} Given a bounded program \(P\), how many Mazurkiewicz equivalence classes do its interleavings induce?
\end{quote}

For an optimal model checking algorithm, the estimation problem is directly related (but not identical) to the cost of model checking with partial order reduction.
Thus, solving the estimation problem can help answer two
high-level quantitative questions that arise when using a model checker:
\emph{``How long should we expect the model checker to run?''} and
\emph{``What fraction of the state space has the model checker covered so far?''}
Reliable answers to these questions allow users of model checking tools to predict verification effort (see, e.g., discussions in \cite{WangCGMK18,gator}).
%
% Reliable answers to these questions have the potential to significantly enhance
% the usability and adoption of model checkers \cite{gator,jpf}, by enabling software teams
% to allocate computational budgets more effectively, predict the remaining
% verification effort, and meaningfully compare competing verification tools.
% As model checkers become an integral part of the modern software development toolkit
% \cite{spin, cbmc,jbmc, p-lang, trust,optimal-dpor,nidhugg,condpor,kani,must,jpf,cpachecker,cdschecker,concuerror,rInspect},
% these resourcing questions become increasingly important.

We first ask whether the above estimation problem can be solved
tractably (i.e., in polynomial time in the size of the program).
The situation is bleak:
we show that the problem is \#P-hard even for highly restricted concurrent programs.
Informally, the \#P-hardness result implies that the task is at least as hard as
counting the satisfying assignments of a Boolean formula (\#SAT), a canonical
counting problem for which no efficient algorithm is believed to exist.
In fact, we show that unless \textsf{P} = \textsf{NP}, we cannot even \emph{approximate} the number of equivalence classes
up to a subexponential factor!
Thus, neither exact nor even moderately precise approximate
counting is feasible in general.

Although these hardness results rule out efficient exact or approximate
counting, they do not prevent us from obtaining statistical or \emph{Monte Carlo}
estimates of the size of the behavioral-space.
The statistical formulation of our problem asks for an \emph{unbiased estimator}:
\begin{quote}
{\bf Unbiased statistical estimation problem:} Given a bounded program \(P\), give a polynomial-time randomized procedure whose expected value is the number of Mazurkiewicz-equivalence classes.
\end{quote}
%

A natural starting point for deriving such an estimator is Knuth's classical
estimator \cite{knuth75}, which approximates the number of leaves in a
combinatorial search tree by performing a randomized exploratory
walk, recording the branching factors encountered along each walk, and
computing a weighted product; this yields an unbiased estimator of the total
number of leaves.


However, the space of Mazurkiewicz traces is a directed acyclic graph (DAG) but not a tree:
different schedules of events can merge into the same equivalence class.
On a DAG, Knuth's procedure is no longer an unbiased estimator.
Leonard Pitt generalized Knuth's procedure so it works for DAGs\cite{pitt}.
While Pitt's estimator is unbiased, its variance is too high to be of practical use: 
even when a program has exactly one maximal trace, the variance can be exponential.

The main main insight we employ in this work that the \trust stateless optimal \dpor algorithm \cite{trust}
can be converted into an unbiased estimator!
This follows because the \trust algorithm enjoys several crucial properties.
First, its optimality guarantee implies that the underlying state space explored by the algorithm defines a tree structure.
Second, since it is stateless, the next decisions depend solely on the current state and not on the history of the search.
Third, each exploration of the algorithm is bounded polynomially in the size of the program.

These observations together imply that the underlying state space traversed by \trust is a tree, whose depth and width are bounded by a quadratic function of the program size,
and whose leaves are exactly the maximal Mazurkiewicz traces.
Thus, estimating the leaves of this tree solves the estimation problem,
and we can apply the Knuth estimator to this (implicitly constructed) tree to obtain an unbiased estimator.
The algorithm is surprisingly simple: we run \trust, but instead of systematically exploring all successors, we randomly pick one successor at each stage.
At the end, we output the product of the number of choices we had along the path.
Note that we do not exhaustively explore all traces; this number can grow exponentially with program size.

While the estimator is unbiased, its variance can still be exponentially large,
reflecting the inherent hardness of approximation.
However, in many practical situations, a large variance often occurs not because of inherent computational complexity but because the underlying tree is ``skewed''.
In the \trust algorithm, some explored paths are linear in the size of the program while others can
be quadratic---a random walk procedure is much more likely to visit the ``shallow'' leaves than the ``deep'' ones.

In order to reduce variance, we use a technique from Monte Carlo estimation called \emph{stochastic enumeration} \cite{rubinstein, vaisman2017} or weighted ensemble \cite{huberkim}.
Instead of one random walk, as in Knuth's estimator, we maintain a small population of partial paths at each depth of the tree.
Whenever a path ends, the population is redistributed to reflect the structure.
The overall stochastic procedure ensures that no single highly unbalanced subtree dominates the
estimate.
This preserves the unbiasedness of estimation while substantially smoothing out the contribution
of deep but large subtrees.
As our evaluation shows, on a diverse set of bounded concurrent programs, the
resulting estimator attains stable estimates with a modest number of samples,
suggesting that this variance-reduction strategy is well suited to realistic
model-checking workloads.

We have implemented our estimator in the \dpor-based stateless Java model checker JMC and evaluated it on a suite of bounded concurrent programs, including concurrent data structures 
and synthetic benchmarks designed to exercise different concurrency patterns.

We show that the estimator based on \trust, combined with stochastic enumeration, consistently achieves stable estimates within a small number of trials---typically a few hundred---even 
for behavior spaces ranging from about \(10^{5}\) to \(10^{6}\)
executions, with relative error usually within a 20\% band.
On the other hand, each aspect of our estimation algorithm remains 
crucial: on these benchmarks, a DAG-based estimator or an estimator that does not perform stochastic enumeration fails to converge.

We note that the combination of \trust, tree-estimation, and stochastic enumeration exploits a number of subtle but important properties of the \trust algorithm.
First, we crucially leverage the fact that the underlying optimal \dpor algorithm is history independent.
In an optimal model checker like Nidhugg~\cite{nidhugg}, every execution graph is visited through a unique path, implicitly inducing a tree.
However, the task of sampling paths in such a tree (in the style of Knuth's estimator) can be non-trivial.
This is because, to decide whether a given execution graph is an extension of another, the model checker may have to consult the history of the exploration, 
 paying exponential cost in the worst case, undermining the poly-time sampling property.
Second, one could try to reduce variance using a number of heuristics that perform importance sampling in various ways.
However, in practice, the performance of importance sampling is very sensitive to good estimates of the relative sizes of different options.
These estimates are difficult to come by for practical benchmarks.
%% other than by computationally expensive nested Monte Carlo searches.

We also explore the theoretical question of whether there is a good approximate counter that uses \emph{subexponential} time.
We find that, surprisingly, there is an unbiased estimator with low variance that runs in time \(\tilde{O}(\sqrt{N})\), where \(N\) is the number of possible maximal traces (here $\tilde{O}$ hides polynomial factors)!
While theoretically optimal (there is an \(\Omega(\sqrt{N})\) lower bound \cite{Stockmeyer85}), the polynomial factors appearing in the algorithm are too large for empirical evaluation.

\myparagraph{Contributions}{
In summary, this paper makes the following contributions: 
\begin{itemize}[topsep=0pt, partopsep=0pt, itemsep=0.25em, parsep=0pt]
  \item  
  We formulate the estimation problem for Mazurkiewicz traces and show that it is \#P-hard 
  and
  inapproximable within any subexponential factor unless \textsf{P} = \textsf{NP}.

  \item 
  We show how the stateless optimal \dpor algorithm \trust can be modified to produce a poly-time unbiased estimation, and show how a population-based
  estimation strategy reduces the variance of the estimation (\cref{trust_sec}, \cref{sec:se}).
  
  \item We evaluate our estimation procedure on a wide variety of shared-memory concurrency benchmarks in the context of the \jmc model checker (\cref{sec:experiment}).
  The summary of our experiments is that our estimator provides accurate and robust estimates of the number of equivalence classes, as well as the cost of model checking,
  while using a small population budget and a modest number of samples.
\end{itemize}
}

%For space reasons, the detailed proofs be found in the full version of the paper~\cite{balasubramanian2026statespaceestimationdporbased}.
